\section{Universal Quantum Computation}
In classical computing, a set of logic gates is considered \textit{universal} 
if any arbitrary Boolean function can be constructed using only gates from that set (e.g., the NAND gate alone is universal). 

In the quantum regime, an arbitrary unitary transformation requires continuous parameters, 
implying an infinite number of possible gates. 

However, the Solovay-Kitaev theorem states that any quantum algorithm can be approximated to arbitrary accuracy using a finite, 
discrete set of quantum gates, with only a polylogarithmic overhead in the number of gates \cite{dawson2005}. 
A common universal gate set used in fault-tolerant quantum computing is the combination of the single-qubit Hadamard and $T$ gates, 
along with the two-qubit CNOT gate: $\{H, T, \text{CNOT}\}$. 